
Müllsammler

Virtuelle Maschinen (VM) wie die Java VM (JVM)
bieten Garbage Collection (GC), um nicht mehr
benötigte (nicht mehr referenzierte) Daten
automatisch aus dem Arbeitsspeicher (RAM) zu entfernen.

Dadurch kann zum Einen blockierter Speicher schnell wieder
freigegeben und zum Anderen ein Programmabsturz durch
die vergessene Freigabe von Speicher vermieden werden.

Zur Umsetzung einer GC gibt es viele Möglichkeiten
(z. B. SmartPointer, Counter Muster QUELLEN).

Eine einfache Möglichkeit ist es, für jeden Variablenwert
einen eigenen Referenzzähler vorzuhalten, also:

- name
- type
- model
- reference

Sobald der Referenzzähler auf Null sänke, würde der
zugehörige Datenwert zerstört und Speicher freigegeben werden.

Natürlich muss sichergestellt sein, dass alle den Wert
manipulierenden Funktionen (z. B. "insert", "remove")
stets den Referenzzähler hoch- bzw. runterzählen.

Die Thematik GC soll in diesem Artikel jedoch nicht
weiter vertieft werden.
